% GENERATED FILE. Edit canonical lessons, then run npm run generate:books.
% canonical-source-hash: fnv1a64:18fad2204cdaba72
% canonical-lessons: JA-C375-shihatsu, JA-C375-shuuden, JA-C375-shuuten, JA-C375-norikosu, JA-C375-norikomu

\chapter{The First Train and the Last}
\label{ch:ja-the-first-train-and-the-last}
\hlchaptermodality{375}

\begin{chapteropening}
\textbf{By the end of this chapter:} \emph{I can say when the first and last trains run.}

Complete the last lesson of chapter 375: I can say when the first and last trains run.
\end{chapteropening}

\section[\texorpdfstring{\ja{しはつ} (shihatsu)}{しはつ (shihatsu)}]{\ja{しはつ} (shihatsu) — the first train of the day}

\label{lesson:JA-C375-shihatsu}



\begin{quote}
Before the new one: say the Japanese for departure (of a train or bus), then the Japanese for to go near.
\end{quote}

\subsection*{You'll want to know: \ja{しはつ}}
\textbf{\ja{しはつ}} — \emph{shihatsu} — \textquotedblleft{}the first train of the day\textquotedblright{}.

\begin{quote}
\textbf{\ja{しはつは} \ja{ごじに} \ja{でます。}} — \emph{shihatsu wa goji ni demasu} — \textquotedblleft{}The first train leaves at five.\textquotedblright{}
\end{quote}

\begin{cousinweb}
\textbf{\ja{しはつ}} (\emph{shihatsu}) is Sino-Japanese (kango): two Chinese characters read with their on'yomi, readings borrowed from Chinese. They are \emph{shi}, \textquotedblleft{}begin\textquotedblright{} (Middle Chinese \emph{syiX}) and \emph{hatsu}, \textquotedblleft{}departure\textquotedblright{} (\emph{pjot}) (KANJIDIC2; Middle Chinese in Baxter's transcription).
\end{cousinweb}

\begin{grammarlens}[title={What you've built}]
One more word for getting around.
\end{grammarlens}

\subsection*{Your turn}
\begin{itemize}
\raggedright
  \item \emph{Say it:} \emph{shihatsu}
  \item \emph{Say it:} \emph{shihatsu}, once more
  \item \emph{Say it:} \emph{shihatsu wa goji ni demasu} — The first train leaves at five.
  \item \emph{Recall:} say the Japanese for departure (of a train or bus), then the Japanese for to go near, then say \emph{shihatsu} again
\end{itemize}

\begin{tcolorbox}[breakable,colback=teal!4,colframe=teal!35!black,title={Before you move on}]
What does \ja{しはつ} mean? (\textquotedblleft{}The first train of the day\textquotedblright{}.) Say it once more.
\end{tcolorbox}
\section[\texorpdfstring{\ja{しゅうでん} (shūden)}{しゅうでん (shūden)}]{\ja{しゅうでん} (shūden) — the last train of the day}

\label{lesson:JA-C375-shuuden}



\begin{quote}
Before the new one: say the Japanese for to go near, then the Japanese for the first train of the day.
\end{quote}

\subsection*{You'll want to know: \ja{しゅうでん}}
\textbf{\ja{しゅうでん}} — \emph{shūden} — \textquotedblleft{}the last train of the day\textquotedblright{}.

\begin{quote}
\textbf{\ja{しゅうでんに} \ja{まにあいました。}} — \emph{shūden ni maniaimashita} — \textquotedblleft{}I made it in time for the last train.\textquotedblright{}
\end{quote}

\begin{cousinweb}
\textbf{\ja{しゅうでん}} (\emph{shūden}) is Sino-Japanese (kango): two Chinese characters read with their on'yomi, readings borrowed from Chinese. They are \emph{shū}, \textquotedblleft{}end\textquotedblright{} (Middle Chinese \emph{tsyuwng}) and \emph{den}, \textquotedblleft{}electricity\textquotedblright{} (\emph{denH}) (KANJIDIC2; Middle Chinese in Baxter's transcription).
\end{cousinweb}

\begin{grammarlens}[title={What you've built}]
One more word for getting around.
\end{grammarlens}

\subsection*{Your turn}
\begin{itemize}
\raggedright
  \item \emph{Say it:} \emph{shūden}
  \item \emph{Say it:} \emph{shūden}, once more
  \item \emph{Say it:} \emph{shūden ni maniaimashita} — I made it in time for the last train.
  \item \emph{Recall:} say the Japanese for to go near, then the Japanese for the first train of the day, then say \emph{shūden} again
\end{itemize}

\begin{tcolorbox}[breakable,colback=teal!4,colframe=teal!35!black,title={Before you move on}]
What does \ja{しゅうでん} mean? (\textquotedblleft{}The last train of the day\textquotedblright{}.) Say it once more.
\end{tcolorbox}
\section[\texorpdfstring{\ja{しゅうてん} (shūten)}{しゅうてん (shūten)}]{\ja{しゅうてん} (shūten) — the terminus; the last stop}

\label{lesson:JA-C375-shuuten}



\begin{quote}
Before the new one: say the Japanese for the first train of the day, then the Japanese for the last train of the day.
\end{quote}

\subsection*{You'll want to know: \ja{しゅうてん}}
\textbf{\ja{しゅうてん}} — \emph{shūten} — \textquotedblleft{}the terminus; the last stop\textquotedblright{}.

\begin{quote}
\textbf{\ja{この} \ja{でんしゃの} \ja{しゅうてんは} \ja{みなとです。}} — \emph{kono densha no shūten wa minato desu} — \textquotedblleft{}This train's last stop is the harbour.\textquotedblright{}
\end{quote}

\begin{cousinweb}
\textbf{\ja{しゅうてん}} (\emph{shūten}) is Sino-Japanese (kango): two Chinese characters read with their on'yomi, readings borrowed from Chinese. They are \emph{shū}, \textquotedblleft{}end\textquotedblright{} (Middle Chinese \emph{tsyuwng}) and \emph{ten}, \textquotedblleft{}point\textquotedblright{} (\emph{temX}) (KANJIDIC2; Middle Chinese in Baxter's transcription).
\end{cousinweb}

\begin{grammarlens}[title={What you've built}]
One more word for getting around.
\end{grammarlens}

\subsection*{Your turn}
\begin{itemize}
\raggedright
  \item \emph{Say it:} \emph{shūten}
  \item \emph{Say it:} \emph{shūten}, once more
  \item \emph{Say it:} \emph{kono densha no shūten wa minato desu} — This train's last stop is the harbour.
  \item \emph{Recall:} say the Japanese for the first train of the day, then the Japanese for the last train of the day, then say \emph{shūten} again
\end{itemize}

\begin{tcolorbox}[breakable,colback=teal!4,colframe=teal!35!black,title={Before you move on}]
What does \ja{しゅうてん} mean? (\textquotedblleft{}The terminus; the last stop\textquotedblright{}.) Say it once more.
\end{tcolorbox}
\section[\texorpdfstring{\ja{のりこす} (norikosu)}{のりこす (norikosu)}]{\ja{のりこす} (norikosu) — to ride past one's stop}

\label{lesson:JA-C375-norikosu}



\begin{quote}
Before the new one: say the Japanese for the last train of the day, then the Japanese for the terminus.
\end{quote}

\subsection*{You'll want to know: \ja{のりこす}}
\textbf{\ja{のりこす}} — \emph{norikosu} — \textquotedblleft{}to ride past one's stop\textquotedblright{}.

\begin{quote}
\textbf{\ja{でんしゃで} \ja{しんぶんを} \ja{よんで、えきを} \ja{のりこしました。}} — \emph{densha de shinbun o yonde, eki o norikoshimashita} — \textquotedblleft{}I read the newspaper on the train and rode past my station.\textquotedblright{}
\end{quote}

Polite \textbf{\ja{のりこします}}, \emph{norikoshimasu}; a godan verb. Past (polite): \textbf{\ja{あにが} \ja{えきを} \ja{のりこしました。}}, \emph{ani ga eki o norikoshimashita}.

\begin{cousinweb}
\textbf{\ja{のりこす}} (\emph{norikosu}) is a native Japanese word (wago), inherited rather than borrowed from Chinese, built from \emph{nori}, \textquotedblleft{}riding\textquotedblright{} (from \emph{noru}) and \emph{kosu}, \textquotedblleft{}to cross over\textquotedblright{}; its kanji were attached later as kun'yomi (JMdict).
\end{cousinweb}

\begin{grammarlens}[title={What you've built}]
One more word for getting around.
\end{grammarlens}

\subsection*{Your turn}
\begin{itemize}
\raggedright
  \item \emph{Say it:} \emph{norikosu}
  \item \emph{Say it:} \emph{norikosu}, once more
  \item \emph{Say it:} \emph{densha de shinbun o yonde, eki o norikoshimashita} — I read the newspaper on the train and rode past my station.
  \item \emph{Recall:} say the Japanese for the last train of the day, then the Japanese for the terminus, then say \emph{norikosu} again
\end{itemize}

\begin{tcolorbox}[breakable,colback=teal!4,colframe=teal!35!black,title={Before you move on}]
What does \ja{のりこす} mean? (\textquotedblleft{}To ride past one's stop\textquotedblright{}.) Say it once more.
\end{tcolorbox}
\section[\texorpdfstring{\ja{のりこむ} (norikomu)}{のりこむ (norikomu)}]{\ja{のりこむ} (norikomu) — to get on board}

\label{lesson:JA-C375-norikomu}



\begin{quote}
Before the new one: say the Japanese for the terminus, then the Japanese for to ride past one's stop.
\end{quote}

\subsection*{You'll want to know: \ja{のりこむ}}
\textbf{\ja{のりこむ}} — \emph{norikomu} — \textquotedblleft{}to get on board\textquotedblright{}.

\begin{quote}
\textbf{\ja{みんなが} \ja{しはつに} \ja{のりこみました。}} — \emph{minna ga shihatsu ni norikomimashita} — \textquotedblleft{}Everyone got on board the first train.\textquotedblright{}
\end{quote}

Polite \textbf{\ja{のりこみます}}, \emph{norikomimasu}; a godan verb. Past (polite): \textbf{\ja{じょうきゃくが} \ja{でんしゃに} \ja{のりこみました。}}, \emph{jōkyaku ga densha ni norikomimashita}.

\begin{cousinweb}
\textbf{\ja{のりこむ}} (\emph{norikomu}) is a native Japanese word (wago), inherited rather than borrowed from Chinese, built from \emph{nori}, \textquotedblleft{}riding\textquotedblright{} (from \emph{noru}) and \emph{komu}, \textquotedblleft{}to go in\textquotedblright{}; its kanji were attached later as kun'yomi (JMdict).
\end{cousinweb}

\begin{grammarlens}[title={What you've built}]
One more word for getting around.
\end{grammarlens}

\subsection*{Your turn}
\begin{itemize}
\raggedright
  \item \emph{Say it:} \emph{norikomu}
  \item \emph{Say it:} \emph{norikomu}, once more
  \item \emph{Say it:} \emph{minna ga shihatsu ni norikomimashita} — Everyone got on board the first train.
  \item \emph{Recall:} say the Japanese for the terminus, then the Japanese for to ride past one's stop, then say \emph{norikomu} again
\end{itemize}

\begin{tcolorbox}[breakable,colback=teal!4,colframe=teal!35!black,title={Before you move on}]
What does \ja{のりこむ} mean? (\textquotedblleft{}To get on board\textquotedblright{}.) Say it once more.
\end{tcolorbox}
